% Front-door note N1 (math.OA audience) — Phase 2a of the corpus
% accessibility plan. NOT SENT; sends are Phase 4, PI-gated.
% Self-contained: no project-internal vocabulary.
\documentclass[11pt]{article}
\usepackage[T1]{fontenc}
\usepackage[utf8]{inputenc}
\usepackage{amsmath, amssymb, amsthm}
\usepackage[margin=1.1in]{geometry}
\usepackage{hyperref}

\newtheorem{theorem}{Theorem}
\newtheorem{lemma}[theorem]{Lemma}
\newtheorem{question}{Question}
\theoremstyle{remark}
\newtheorem{remark}{Remark}

\newcommand{\R}{\mathbb{R}}
\newcommand{\C}{\mathbb{C}}
\newcommand{\N}{\mathbb{N}}
\newcommand{\Op}{\mathcal{O}}
\newcommand{\opnorm}[1]{\Vert #1 \Vert}
\newcommand{\norm}[1]{\Vert #1 \Vert}
\newcommand{\DCH}{D}

\title{Spectral truncations of the Dirac triple on $S^3$:\\
a convergence theorem at explicit rate, a degeneracy obstruction\\
for the obvious Lorentzian extension, and three questions}
\author{J.~Loutey\thanks{Independent researcher,
\texttt{jloutey@gmail.com}. This note summarizes two write-ups
developed in an AI-assisted workflow (implementation and drafting
with Anthropic's Claude;\ all claims and their verification status
are stated below exactly as established, and all errors are the
author's). Full papers, code, and bit-exact verification scripts:\
\texttt{https://github.com/jloutey-hash/geovac}.}}
\date{June 2026}

\begin{document}
\maketitle

\noindent This is a three-page summary of two results on
Connes--van~Suijlekom spectral truncations
\cite{CvS2021}, written to ask three concrete questions of
specialists. Verification status is flagged sentence by
sentence;\ proofs are self-contained in the full write-ups, and
anything verified only numerically is flagged as such.

\section{Setting}

Let $\mathcal{T}_{S^3} = (C^\infty(S^3), L^2(S^3, \Sigma), \DCH)$ be
the spectral triple of the Dirac operator on the round unit
$S^3 = \mathrm{SU}(2)$, with spectrum $\pm(n + \tfrac32)$,
$n \ge 0$, and multiplicity $(n+1)(n+2)$ per sign. For a cutoff
$\Lambda$, let $P_\Lambda$ be the spectral projection and
$\Op_\Lambda = P_\Lambda C^\infty(S^3) P_\Lambda$ the truncated
operator system \cite{CvS2021}. Van~Suijlekom \cite{vS2021} proved
Gromov--Hausdorff convergence of the truncated state spaces
$S(\Op_\Lambda) \to S(C(S^1))$ for the circle, and
Gaudillot-Estrada--van~Suijlekom \cite{GEvS2023} extended state-space
GH convergence to scalar truncations of all compact metric groups ---
including $\mathrm{SU}(2)$. What the present work adds on the
Riemannian side is therefore \emph{not} existence of convergence but:\
the spinor-bundle (Dirac) substrate, an explicit closed-form rate
moment, and the asymptotic rate constant. On the Lorentzian side it
adds a negative structural theorem.

\section{Construction and rate data}

On $\mathrm{SU}(2)$, define the central spectral Fej\'er kernel
\begin{equation}
  K_n(g) \;=\; \frac{1}{Z_n}\Bigl|\sum_{j \le (n-1)/2}
  \sqrt{2j+1}\,\chi_j(g)\Bigr|^2, \qquad Z_n = \tfrac{n(n+1)}{2},
\end{equation}
a positive, central, unit-mass kernel ($\chi_j$ = trace characters).
Its smoothing map $f \mapsto K_n * f$ is unital completely positive
(cb-norm $1$);\ its multiplier symbol is given by a Clebsch--Gordan
sum with $\sigma(0) = 1$, $0 \le \sigma \le 1$, support
$j' \le n - 1$. Its mass-concentration moment
\begin{equation}
  \gamma_n \;=\; \int_{\mathrm{SU}(2)} K_n(g)\,
  d_{\mathrm{round}}(e, g)\,dg
  \;=\; \pi - \frac{4 T_n}{\pi Z_n}, \qquad
  T_n = \!\!\sum_{\substack{1 \le k_1, k_2 \le n \\ k_1 + k_2\
  \mathrm{odd}}}\!\! \sqrt{k_1 k_2}\Bigl[\frac{1}{(k_1 - k_2)^2}
  - \frac{1}{(k_1 + k_2)^2}\Bigr],
\end{equation}
satisfies $\gamma_n \sim \frac{4}{\pi}\,\frac{\log n}{n}$ --- the
same constant $4/\pi$ as the classical Fej\'er--Stein--Weiss circle
estimate, here arising as
$\mathrm{Vol}(S^2)/\pi^2$ on the Hopf quotient
$\mathrm{SU}(2)/U(1)$. A Lipschitz comparison
$\opnorm{[\DCH, M_f]} \le C_3 \norm{\nabla f}_\infty$ holds with
$C_3 = 1$ (per-harmonic constant $(N-1)/\sqrt{N^2-1}$), and the
convolution-form Berezin map
$B_n(f) = P_n M_{K_n * f} P_n$ is UCP, an approximate identity at
rate $\gamma_n$, and Lipschitz-compatible.

\section{Result A:\ unconditional convergence}

\begin{theorem}
The truncations of the chirality-doubled Dirac triple above converge
to $\mathcal{T}_{S^3}$ in van~Suijlekom's state-space
Gromov--Hausdorff distance, with the truncated state spaces metrized
by the left-translation Lipschitz seminorm
$L_n(T) = \sup_{g \ne e} \norm{U_g T U_g^* - T} / d(e, g)$ (equal to
the metric Lipschitz constant on the continuum):
$d_{\mathrm{GH}} \le \gamma_n = \tfrac{4 \log n}{\pi n} + O(1/n)$.
\end{theorem}

\noindent\emph{Proof shape} (self-contained in the full write-up).
The kernel condition $L_n(T) = 0 \Leftrightarrow T \in \C 1$ holds on
the truthful substrate:\ left translation preserves harmonic bands,
band-limited compression is injective band-by-band (Schur
multiplicity-one;\ the nonvanishing matrix element is the stretched
$3$-$Y$ coefficient), and transitivity does the rest. The dual map
evaluates conjugated operators against the Fej\'er vector state
$\xi = (h \otimes \chi)/\sqrt{Z}$, which lies \emph{exactly} in the
Dirac window:\ under $\mathrm{SU}(2)_L \times \mathrm{SU}(2)_R$ the
spinor space is $\bigoplus_j V_j \otimes (V_{j+1/2} \oplus
V_{j-1/2})$ and these summands are precisely the Dirac eigenspaces.
Both almost-inverse defects collapse to Fej\'er smoothing — one on
functions, one as the operator conjugation-average — at the common
moment $\gamma_n$, so no partial-inverse or transference estimate
appears anywhere. The truthful Dirac-\emph{commutator} seminorm, by
contrast, degenerates at finite cutoff (kernel $10/14$ at $n = 2$,
$26/55$ at $n = 3$, bit-exact) — it is the wrong finite-cutoff
metrization, and this degeneracy is quantified in the write-up. The
scalar $C(G)$ pattern is \cite{GEvS2023};\ the contribution here is
the spinor-window transport, the per-band injectivity, the
closed-form rate moment, and the $4/\pi$ constant.

\section{Result B:\ a degeneracy theorem for the obvious Lorentzian
extension}

Take the compact-temporal product $S^3 \times S^1_T$ with Krein space
$\mathcal{K} = \mathcal{H} \otimes \C^{N_t}$, fundamental symmetry
$J = J_s \otimes I$, Lorentzian Dirac
$D_L = i(\gamma^0 \otimes \partial_t + \DCH \otimes I)$
(Krein-self-adjoint, van~den~Dungen's convention \cite{vdD2016}),
and the natural operator system of spatial multipliers tensored with
momentum-diagonal temporal multipliers. The obvious route to a
``Lorentzian quantum metric'' is to compress to the Krein-positive
eigenspace $\mathcal{K}^+$ of $J$ (where the Krein product is a
Hilbert inner product) and run Riemannian machinery on
$(P_+ \Op P_+,\, \mathcal{K}^+,\, P_+ D_L P_+)$.

\begin{theorem}[unconditional]
This route is structurally closed. Krein-self-adjointness of the
anti-Hermitian spatial block forces
$\{J, \DCH \otimes I\} = 0$, hence
$P_+ (\DCH \otimes I) P_+ = 0$ exactly:\ the compression annihilates
the spatial Dirac. The surviving temporal block is momentum-diagonal
and commutes with every momentum-diagonal multiplier. Consequently
the compressed-Dirac Lipschitz seminorm vanishes \emph{identically}
on the operator system, and no quantum Gromov--Hausdorff-type
distance is defined.
\end{theorem}

\begin{proof}[Proof sketch]
$(iD_s)^\dagger = -iD_s$ with $D_s := \DCH \otimes I$, so
$J D_L^\dagger J = D_L$ forces $J D_s J = -D_s$;\ then
$P_+ D_s P_+ = \tfrac14(D_s + D_s J + J D_s + J D_s J) =
\tfrac14(D_s + 0 - D_s) = 0$. The temporal commutators vanish since
both factors are Fourier-diagonal. Verified bit-exactly at two panel
cells (seminorm $= 0.0$ in \texttt{float64} on $42/42$ and $275/275$
basis multipliers;\ a frozen regression test guards the statement).
\end{proof}

The failure analysis sharpens the open problem:\ any viable
Lorentzian/Krein convergence construction needs (i) a temporal
multiplier algebra the Dirac can see --- Toeplitz compressions of
$e^{iqt}$ in the Connes--van~Suijlekom circle pattern, not
momentum-diagonal functions, whose commutators with $\partial_t$
vanish identically;\ (ii) a Krein-compatible Lipschitz seminorm that
does not compress $D_L$ (candidates:\ seminorms built from $JD_L$ or
$|D_L|$, or a genuinely indefinite Lipschitz analog);\ and (iii) a
substrate on which $L(a) = 0 \Leftrightarrow a \in \C 1$ actually
holds.

\section{Three questions}

\begin{question}
Is the degeneracy theorem (Result B) known --- in any form --- in the
Krein-spectral-triple or Lorentzian NCG literature? It is elementary
once stated, which makes prior art likely;\ we could not locate it.
\end{question}

\begin{question}
For Result A:\ is the finite-cutoff degeneracy of the
Dirac-commutator seminorm — and its repair by the translation
seminorm — known or recorded? And is the exact-fit spinor-window
embedding of the Fej\'er vector state (via
$V_j \otimes V_{j \pm 1/2}$) the construction you would have used,
or is there a slicker route to the dual map on bundle-valued
truncations?
\end{question}

\begin{question}
Is the Lorentzian convergence problem, as specified by (i)--(iii)
above, of interest to the quantum-metric community --- and is anyone
already working on a temporal-Toeplitz version?
\end{question}

\begin{thebibliography}{9}
\bibitem{CvS2021} A.~Connes, W.~D.~van~Suijlekom,
\emph{Spectral truncations in noncommutative geometry and operator
systems}, Comm. Math. Phys. \textbf{383} (2021) 2021--2067.
\bibitem{vS2021} W.~D.~van~Suijlekom,
\emph{Gromov--Hausdorff convergence of state spaces for spectral
truncations}, J. Geom. Phys. \textbf{162} (2021) 104075.
\bibitem{GEvS2023} Y.~Gaudillot-Estrada, W.~D.~van~Suijlekom,
\emph{Convergence of spectral truncations for compact metric groups},
arXiv:2310.14733.
\bibitem{LvS2024} M.~Leimbach, W.~D.~van~Suijlekom,
\emph{Gromov--Hausdorff convergence of spectral truncations of the
torus}, Adv. Math. \textbf{439} (2024) 109496.
\bibitem{vdD2016} K.~van~den~Dungen,
\emph{Krein spectral triples and the fermionic action},
Math. Phys. Anal. Geom. \textbf{19} (2016), Art.~4.
\end{thebibliography}

\end{document}
